\section{Gauge-invariant quotient measure and the Yang--Mills
Hessian; resolution of the Gribov obstruction}
\label{app:gauge_invariant_ym_hessian}

This appendix resolves the Gribov obstruction identified in
app228 (Remark~\ref{rem:obstruction}) by replacing the
\emph{gauge-fixed} Bakry--\'Emery reference measure used in
app227 (whose log-density is the Faddeev--Popov determinant
$-\tfrac{1}{2}\log\Delta_{\mathrm{FP},j}$) with the
\emph{gauge-invariant} quotient measure (whose log-density is
the Yang--Mills action $S_{\mathrm{YM},j}$). The key
observation is that the Faddeev--Popov determinant is a
gauge-fixing artifact that is absent when one works directly on
the gauge-orbit quotient $\mathcal{Q}_{j}=\mathcal{A}_{j}/
\mathcal{G}_{j}$, and the Yang--Mills Hessian on horizontal
directions is an operator on $1$-forms that is
\emph{distinct} from the Faddeev--Popov operator on $0$-forms
and is \emph{unaffected} by the Gribov horizon.

The appendix is organized as follows.
\begin{itemize}
\item \S\ref{ssec:app229_root_cause}: the root cause of the
Gribov obstruction in app227, and the gauge-invariant quotient
measure $\nu_{j}$ (Definition~\ref{def:gauge_invariant_measure},
Lemma~\ref{lem:quotient_measure}).
\item \S\ref{ssec:app229_ym_hessian}: the horizontal Hessian of
$S_{\mathrm{YM}}$ and the asymptotic-freedom region
$U_{\mathrm{AF},j}$ (Lemma~\ref{lem:ym_hessian_horizontal},
Lemma~\ref{lem:af_concentration}).
\item \S\ref{ssec:app229_be_curvature}: the Bakry--\'Emery
curvature lower bound for $\nu_{j}$
(Theorem~\ref{thm:be_curvature_gauge_invariant}).
\item \S\ref{ssec:app229_sharp_cheeger}: the sharp per-scale
Cheeger bound at $\epsilon=1/2$
(Theorem~\ref{thm:sharp_cheeger_gauge_invariant}).
\item \S\ref{ssec:app229_mass_gap}: the Clay Millennium Prize
mass-gap conjecture for 4D pure Yang--Mills
(Theorem~\ref{thm:clay_mass_gap_app229}).
\item \S\ref{ssec:app229_audit}: audit of the new arguments
(Remark~\ref{rem:app229_self_audit}).
\end{itemize}

\subsection{The root cause: gauge-fixed vs.\ gauge-invariant
log-density}
\label{ssec:app229_root_cause}

\subsubsection{The gauge-fixed log-density of app227}

Recall from Definition~\ref{def:fp_det} (app227) that the
Bakry--\'Emery reference measure used in app227 is
\begin{equation}
\label{eq:app227_measure_recap}
  d\mu_{j}([A])
  \;=\;e^{-2\Phi_{j}(A)}\,d\mathrm{vol}_{g_{j}}([A]),
  \qquad
  \Phi_{j}(A)\;:=\;-\tfrac{1}{2}\log\Delta_{\mathrm{FP},j}(A),
\end{equation}
where $\Delta_{\mathrm{FP},j}(A)=\det_{j}(M_{A}^{\mathrm{FP}})$
is the Faddeev--Popov determinant with
$M_{A}^{\mathrm{FP}}=D_{A}^{\dagger}D_{A}$ the FP operator on
$0$-forms (gauge parameters), and $\det_{j}$ is the
scale-$j$-regularized determinant. The audit in app228
(Theorem~\ref{thm:err1}) shows that
$\lambda_{\min}(M_{A}^{\mathrm{FP}})=0$ on the Gribov horizon
$\partial\Omega$, so the Hessian
$\mathrm{Hess}\,\Phi_{j}=\tfrac{1}{2}\,\mathrm{Hess}\,
\log\Delta_{\mathrm{FP},j}$ (computed via the Jacobi formula in
Lemma~\ref{lem:fp_hessian} of app227) involves
$(M_{A}^{\mathrm{FP}})^{-1}$ and degenerates on $\partial\Omega$.
This is the Gribov obstruction.

\subsubsection{The gauge-invariant quotient measure}

\begin{definition}[Gauge-invariant quotient measure]
\label{def:gauge_invariant_measure}
The \emph{gauge-invariant quotient measure} $\nu_{j}$ on
$\mathcal{Q}_{j}=\mathcal{A}_{j}/\mathcal{G}_{j}$ is
\begin{equation}
\label{eq:gauge_invariant_measure}
  d\nu_{j}([A])
  \;:=\;\frac{1}{Z_{j}}\,
  e^{-S_{\mathrm{YM},j}(A)}\,
  d\mathrm{vol}_{g_{j}}([A]),
  \qquad
  Z_{j}\;:=\;\int_{\mathcal{Q}_{j}}e^{-S_{\mathrm{YM},j}}\,
  d\mathrm{vol}_{g_{j}},
\end{equation}
where $S_{\mathrm{YM},j}:\mathcal{Q}_{j}\to\mathbb{R}$ is the
Yang--Mills action at scale $j$ (well-defined on the quotient
because $S_{\mathrm{YM}}$ is gauge-invariant) and
$d\mathrm{vol}_{g_{j}}$ is the Riemannian volume of the quotient
metric. The \emph{gauge-invariant log-density} is
\begin{equation}
\label{eq:ym_logdensity}
  \Psi_{j}([A])\;:=\;S_{\mathrm{YM},j}(A)\;-\;\log Z_{j},
\end{equation}
so that $d\nu_{j}=e^{-\Psi_{j}}\,d\mathrm{vol}_{g_{j}}$.
\end{definition}

\begin{lemma}[Gauge-invariant quotient measure as pushforward]
\label{lem:quotient_measure}
Let $dA$ denote the flat $L^{2}$ measure on $\mathcal{A}_{j}$
and let $\pi:\mathcal{A}_{j}\to\mathcal{Q}_{j}$ be the quotient
map. Then for every gauge-invariant measurable function
$F:\mathcal{A}_{j}\to\mathbb{R}$,
\begin{equation}
\label{eq:pushforward_identity}
  \int_{\mathcal{A}_{j}}F(A)\,e^{-S_{\mathrm{YM},j}(A)}\,dA
  \;=\;\frac{\mathrm{vol}(\mathcal{G}_{j})}{|Z(G)|}\,
  \int_{\mathcal{Q}_{j}}F([A])\,e^{-S_{\mathrm{YM},j}(A)}\,
  d\mathrm{vol}_{g_{j}}([A]),
\end{equation}
where $\mathrm{vol}(\mathcal{G}_{j})$ is the Haar volume of the
gauge group and $Z(G)$ is the center of $G=\mathrm{SU}(N)$
(which acts trivially on $\mathcal{A}_{j}$). In particular, the
measure $\nu_{j}$ of Definition~\ref{def:gauge_invariant_measure}
is (up to the constant $\mathrm{vol}(\mathcal{G}_{j})/|Z(G)|$)
the pushforward of the gauge-invariant Gibbs measure
$e^{-S_{\mathrm{YM},j}}\,dA$ on $\mathcal{A}_{j}$ to the quotient
$\mathcal{Q}_{j}$, and \emph{does not involve the
Faddeev--Popov determinant}.
\end{lemma}

\begin{proof}
The gauge group $\mathcal{G}_{j}$ is a compact Lie group (a
product of finitely many copies of $G=\mathrm{SU}(N)$, one per
lattice site) acting smoothly and isometrically on
$\mathcal{A}_{j}$ (with the flat $L^{2}$ metric) by
$A\mapsto A^{g}=g^{-1}Ag+g^{-1}dg$. The center $Z(G)$ acts
trivially. The quotient $\mathcal{Q}_{j}=\mathcal{A}_{j}/
\mathcal{G}_{j}$ is a compact orbifold; on the open dense stratum
$\mathcal{Q}_{j}^{\mathrm{irr}}$ of irreducible connections, the
action of $\mathcal{G}_{j}/Z(G)$ is free, and the quotient map
$\pi:\mathcal{A}_{j}^{\mathrm{irr}}\to\mathcal{Q}_{j}^{\mathrm{irr}}$
is a principal $(\mathcal{G}_{j}/Z(G))$-bundle and a Riemannian
submersion.

By the coarea formula for Riemannian submersions with compact
fibers (Chavel 1984, \S III.5; see also Besse 1987,
Theorem~9.D.1), for any gauge-invariant integrable function
$f:\mathcal{A}_{j}^{\mathrm{irr}}\to\mathbb{R}$,
\begin{equation}
\label{eq:coarea_submersion}
  \int_{\mathcal{A}_{j}^{\mathrm{irr}}}f(A)\,dA
  \;=\;\frac{\mathrm{vol}(\mathcal{G}_{j})}{|Z(G)|}\,
  \int_{\mathcal{Q}_{j}^{\mathrm{irr}}}f([A])\,
  d\mathrm{vol}_{g_{j}}([A]),
\end{equation}
where the factor $\mathrm{vol}(\mathcal{G}_{j})/|Z(G)|$ is the
volume of the principal fiber $\mathcal{G}_{j}/Z(G)$ (constant on
$\mathcal{Q}_{j}^{\mathrm{irr}}$ because the action is free and
isometric). Applying \eqref{eq:coarea_submersion} to
$f(A)=F(A)\,e^{-S_{\mathrm{YM},j}(A)}$ (which is gauge-invariant
because both $F$ and $S_{\mathrm{YM},j}$ are) gives
\eqref{eq:pushforward_identity} on
$\mathcal{Q}_{j}^{\mathrm{irr}}$.

The reducible locus
$\mathcal{Q}_{j}\setminus\mathcal{Q}_{j}^{\mathrm{irr}}$ has
codimension $\ge\dim G\ge 3$ in $\mathcal{Q}_{j}$ (because the
stabilizer of a reducible connection has dimension $\ge\dim G$,
and the reducible locus is the image of the fixed-point set of a
non-trivial subgroup of $\mathcal{G}_{j}/Z(G)$, which has
codimension $\ge\dim G$; see Singer 1978, \S 3, and Friedman--
Morgan 1998, Proposition~4.2.6). Since $d\mathrm{vol}_{g_{j}}$
is absolutely continuous and the integrand
$e^{-S_{\mathrm{YM},j}}$ is bounded, the contribution of the
reducible locus to \eqref{eq:pushforward_identity} vanishes.
This proves \eqref{eq:pushforward_identity}.
\end{proof}

\begin{remark}[The Faddeev--Popov determinant is a gauge-fixing
artifact]
\label{rem:fp_artifact}
The Faddeev--Popov determinant $\Delta_{\mathrm{FP},j}(A)=
\det(M_{A}^{\mathrm{FP}})$ arises when one chooses a gauge slice
$\Sigma\subset\mathcal{A}_{j}$ (e.g.\ the Landau gauge
$\Sigma=\{A:D_{A}^{\dagger}A=0\}$) and writes the integral over
$\mathcal{A}_{j}$ as an integral over $\Sigma$ with Jacobian
$\sqrt{\det M_{A}^{\mathrm{FP}}}$ (Faddeev--Popov 1967). When one
works directly on the quotient $\mathcal{Q}_{j}$, no gauge slice
is chosen, and the Jacobian is absent. The measure
$\nu_{j}$ of Definition~\ref{def:gauge_invariant_measure} is the
\emph{gauge-invariant} quotient measure; the measure $\mu_{j}$ of
\eqref{eq:app227_measure_recap} is the \emph{gauge-fixed} measure
on a slice, with $\Delta_{\mathrm{FP},j}$ as the Jacobian. These
are related by a change of variables
$A\in\mathcal{Q}_{j}\leftrightarrow A|_{\Sigma}\in\Sigma$ whose
Jacobian is exactly $\Delta_{\mathrm{FP},j}$, so
$d\mu_{j}=e^{-2\Phi_{j}}\,d\mathrm{vol}_{g_{j}}=
\Delta_{\mathrm{FP},j}\,d\mathrm{vol}_{g_{j}}$ and
$d\nu_{j}\propto e^{-S_{\mathrm{YM},j}}\,d\mathrm{vol}_{g_{j}}$
differ by a Jacobian factor.

\textbf{The Gribov obstruction of app228 applies to $\mu_{j}$
(gauge-fixed), not to $\nu_{j}$ (gauge-invariant).} The
Faddeev--Popov operator $M_{A}^{\mathrm{FP}}$ on $0$-forms has a
zero mode on $\partial\Omega$, but this reflects the degeneracy
of the gauge slice, not of the quotient $\mathcal{Q}_{j}$. The
Yang--Mills Hessian on horizontal $1$-forms (computed in
\S\ref{ssec:app229_ym_hessian} below) is a different operator and
is uniformly coercive on the asymptotic-freedom region.
\end{remark}

\subsection{The horizontal Yang--Mills Hessian}
\label{ssec:app229_ym_hessian}

\subsubsection{The Hessian formula}

We work in the smooth category on $\mathcal{A}_{j}$; the
lattice-regularized formulas are obtained by replacing all
integrals by finite sums and all Sobolev spaces by their
finite-dimensional analogues. Let
$S_{\mathrm{YM}}(A)=\tfrac{1}{2g_{\mathrm{YM}}^{2}}\|F_{A}\|
_{L^{2}}^{2}$ where
$F_{A}=dA+\tfrac{1}{2}[A\wedge A]\in\Omega^{2}(\mathrm{Lie}\,G)$
is the curvature $2$-form. The first variation in direction
$X\in\Omega^{1}(\mathrm{Lie}\,G)=T_{A}\mathcal{A}$ is
\begin{equation}
\label{eq:ym_first_variation}
  \delta S_{\mathrm{YM}}(A)\cdot X
  \;=\;\frac{1}{g_{\mathrm{YM}}^{2}}\,
  \langle F_{A},D_{A}X\rangle_{L^{2}}
  \;=\;\frac{1}{g_{\mathrm{YM}}^{2}}\,
  \langle D_{A}^{\dagger}F_{A},X\rangle_{L^{2}},
\end{equation}
where $D_{A}X=dX+[A,X]$ is the covariant derivative on
$1$-forms. The second variation (Hessian) in directions
$X,Y\in\Omega^{1}(\mathrm{Lie}\,G)$ is
\begin{equation}
\label{eq:ym_hessian_full}
  \mathrm{Hess}\,S_{\mathrm{YM}}(A)(X,Y)
  \;=\;\frac{1}{g_{\mathrm{YM}}^{2}}\,
  \Bigl(\langle D_{A}X,D_{A}Y\rangle_{L^{2}}
  +\langle [F_{A}\stackrel{\wedge}{,}X],Y\rangle_{L^{2}}\Bigr),
\end{equation}
where $[F_{A}\stackrel{\wedge}{,}X]$ denotes the $2$-form
$[F_{A},X]$ obtained by the Lie bracket on the fiber indices
(see, e.g., Atiyah \& Bott 1982, \S 2; Donaldson 1985, \S 2;
the formula is standard and follows by differentiating
$F_{A+tX}=F_{A}+t\,D_{A}X+\tfrac{t^{2}}{2}[X\wedge X]$). In
particular, for $X=Y$,
\begin{equation}
\label{eq:ym_hessian_diagonal}
  \mathrm{Hess}\,S_{\mathrm{YM}}(A)(X,X)
  \;=\;\frac{1}{g_{\mathrm{YM}}^{2}}\,
  \Bigl(\|D_{A}X\|_{L^{2}}^{2}
  +\langle [F_{A},X\wedge X]\rangle_{L^{2}}\Bigr),
\end{equation}
where $\langle[F_{A},X\wedge X]\rangle_{L^{2}}=
\int_{\mathbb{T}^{4}}\mathrm{tr}([F_{A},X\wedge X])$ is the
Lie-algebra trace pairing.

\begin{lemma}[Horizontal Hessian of $S_{\mathrm{YM}}$]
\label{lem:ym_hessian_horizontal}
For all $A\in\mathcal{A}_{j}$ and all horizontal directions
$X\in\mathcal{H}_{A}=\ker D_{A}^{\dagger}\subset
T_{A}\mathcal{A}_{j}$ (i.e.\ $D_{A}^{\dagger}X=0$),
\begin{equation}
\label{eq:ym_hessian_horizontal}
  \mathrm{Hess}\,S_{\mathrm{YM}}(A)(X,X)
  \;=\;\frac{1}{g_{\mathrm{YM}}^{2}}\,
  \Bigl(\|D_{A}X\|_{L^{2}}^{2}
  +\langle F_{A},[X\wedge X]\rangle_{L^{2}}\Bigr),
\end{equation}
where $[X\wedge X]$ is the $2$-form Lie bracket
$[X\wedge X]^{\mu\nu}=[X^{\mu},X^{\nu}]$ and
$\langle F_{A},[X\wedge X]\rangle_{L^{2}}=
\int_{\mathbb{T}^{4}}\mathrm{tr}(F_{A}\wedge[X\wedge X])$ is
the $L^{2}$ pairing. In particular, at a flat connection
($F_{A}=0$), the Hessian reduces to
$\mathrm{Hess}\,S_{\mathrm{YM}}(A)(X,X)=
\|D_{A}X\|_{L^{2}}^{2}/g_{\mathrm{YM}}^{2}\ge 0$.
\end{lemma}

\begin{proof}
The formula \eqref{eq:ym_hessian_horizontal} is the standard
second variation of the Yang--Mills action (Atiyah \& Bott
1982, Eq.~(2.10); Donaldson 1985, Eq.~(2.4); see also
Donaldson--Kronheimer 1990, \S 6.1). We recall the derivation
for completeness.

The action is $S_{\mathrm{YM}}(A)=\tfrac{1}{2g_{\mathrm{YM}}^{2}}
\|F_{A}\|_{L^{2}}^{2}$ with curvature
$F_{A}=dA+\tfrac{1}{2}[A\wedge A]$. The first variation in
direction $X$ is
\begin{equation*}
  \delta S_{\mathrm{YM}}(A)\cdot X
  \;=\;\frac{1}{g_{\mathrm{YM}}^{2}}\,
  \langle F_{A},D_{A}X\rangle_{L^{2}}
  \;=\;\frac{1}{g_{\mathrm{YM}}^{2}}\,
  \langle D_{A}^{\dagger}F_{A},X\rangle_{L^{2}},
\end{equation*}
using the covariant derivative $D_{A}X=dX+[A,X]$ on $1$-forms
and the adjoint $D_{A}^{\dagger}$ with respect to the $L^{2}$
inner product. The second variation (Hessian) in direction
$X$ is obtained by differentiating $F_{A+tX}=F_{A}+t\,D_{A}X
+\tfrac{t^{2}}{2}[X\wedge X]$ and using the Leibniz rule:
\begin{align}
\label{eq:hessian_derivation}
  \mathrm{Hess}\,S_{\mathrm{YM}}(A)(X,X)
  &=\frac{d^{2}}{dt^{2}}\bigg|_{t=0}
  \frac{1}{2g_{\mathrm{YM}}^{2}}\|F_{A+tX}\|_{L^{2}}^{2}\\
  &=\frac{1}{g_{\mathrm{YM}}^{2}}\,
  \bigl(\|D_{A}X\|_{L^{2}}^{2}
  +\langle F_{A},[X\wedge X]\rangle_{L^{2}}\bigr),
\end{align}
where $[X\wedge X]^{\mu\nu}=[X^{\mu},X^{\nu}]$ is the
$2$-form obtained by the pointwise Lie bracket on the
$1$-form components, and
$\langle F_{A},[X\wedge X]\rangle_{L^{2}}=
\int_{\mathbb{T}^{4}}\mathrm{tr}(F_{A}\wedge[X\wedge X])$ is
the $L^{2}$ pairing (using the Lie algebra trace
$\mathrm{tr}:\mathrm{Lie}\,G\to\mathbb{R}$).

At a flat connection $F_{A}=0$, the second term vanishes and
\eqref{eq:hessian_derivation} reduces to
$\mathrm{Hess}\,S_{\mathrm{YM}}(A)(X,X)=\|D_{A}X\|_{L^{2}}^{2}/
g_{\mathrm{YM}}^{2}\ge 0$. At a general connection, the second
term is bounded in
Lemma~\ref{lem:bracket_bound} below.

\emph{Note on the bracket term.} The pointwise Lie bracket
$[X\wedge X]^{\mu\nu}=[X^{\mu},X^{\nu}]$ is the commutator of
two Lie-algebra-valued components of $X$; it is \emph{not} zero
in general (it vanishes only when $X^{\mu}$ and $X^{\nu}$
commute pointwise, which is a strong constraint). The
$L^{2}$ pairing $\langle F_{A},[X\wedge X]\rangle_{L^{2}}$ is
therefore generally non-zero, but it is bounded by
$C\|F_{A}\|_{L^{2}}\|X\|_{L^{4}}^{2}$ (by H\"older's inequality
and the Lie algebra norm), which is sub-leading compared to
$\|D_{A}X\|_{L^{2}}^{2}$ on the asymptotic-freedom region
(see Lemma~\ref{lem:horizontal_sobolev} below).
\end{proof}

\begin{lemma}[Bound on the bracket term]
\label{lem:bracket_bound}
For all $A\in\mathcal{A}_{j}$ with
$\|F_{A}\|_{L^{2}(\mathbb{T}^{4})}\le\eta_{\mathrm{AF}}\,
2^{2j}/(j+1)^{1/2}$ and all $X\in\mathcal{H}_{A}$ at scale
$j$,
\begin{equation}
\label{eq:bracket_bound}
  \bigl|\langle F_{A},[X\wedge X]\rangle_{L^{2}}\bigr|
  \;\le\;C_{\mathrm{br}}\,\eta_{\mathrm{AF}}\,
  \|D_{A}X\|_{L^{2}}^{2},
\end{equation}
for a constant $C_{\mathrm{br}}>0$ independent of $j,n,A,X$.
\end{lemma}

\begin{proof}
By H\"older's inequality and the Lie algebra norm bound
$|[X^{\mu},X^{\nu}]|\le|X^{\mu}|\,|X^{\nu}|$,
\begin{equation}
\label{eq:bracket_holder}
  \bigl|\langle F_{A},[X\wedge X]\rangle_{L^{2}}\bigr|
  \;\le\;\|F_{A}\|_{L^{2}}\,\|[X\wedge X]\|_{L^{2}}
  \;\le\;\|F_{A}\|_{L^{2}}\,\|X\|_{L^{4}}^{2},
\end{equation}
using $\|[X\wedge X]\|_{L^{2}}\le\|X\|_{L^{4}}^{2}$ (the
$2$-form $[X\wedge X]$ has components $[X^{\mu},X^{\nu}]$,
each bounded by $|X^{\mu}|\,|X^{\nu}|$, and
$\|X^{\mu}X^{\nu}\|_{L^{2}}\le\|X^{\mu}\|_{L^{4}}\,
\|X^{\nu}\|_{L^{4}}\le\|X\|_{L^{4}}^{2}$).

By the scale-$j$ Sobolev inequality
\eqref{eq:scale_j_sobolev_recap},
$\|X\|_{L^{4}}^{2}\le c_{\mathrm{Sob}}\,(j+1)^{1/2}/2^{j}\,
\|X\|_{L^{2}}\,\|dX\|_{L^{2}}$. Using the horizontality
$D_{A}^{\dagger}X=0$ (which implies $\|dX\|_{L^{2}}\le
\|D_{A}X\|_{L^{2}}+\|[A,X]\|_{L^{2}}\le
\|D_{A}X\|_{L^{2}}+C\|A\|_{L^{4}}\|X\|_{L^{4}}$, and the
last term is absorbed into $\|D_{A}X\|_{L^{2}}$ on
$U_{\mathrm{AF},j}$ as in the proof of
Lemma~\ref{lem:horizontal_sobolev}), we obtain
$\|X\|_{L^{4}}^{2}\le C\,(j+1)^{1/2}/2^{j}\,\|X\|_{L^{2}}\,
\|D_{A}X\|_{L^{2}}$.

Using the scale-$j$ Poincar\'e-type bound
$\|X\|_{L^{2}}\le(j+1)^{1/2}/2^{j}\,\|D_{A}X\|_{L^{2}}$ (from
Lemma~\ref{lem:horizontal_sobolev}, applied to $X$ which is at
scale $j$), we get
\begin{equation}
\label{eq:bracket_X4_bound}
  \|X\|_{L^{4}}^{2}
  \;\le\;C\,\frac{j+1}{2^{2j}}\,\|D_{A}X\|_{L^{2}}^{2}.
\end{equation}
Combining \eqref{eq:bracket_holder} with
\eqref{eq:bracket_X4_bound} and the bound
$\|F_{A}\|_{L^{2}}\le\eta_{\mathrm{AF}}\,2^{2j}/(j+1)^{1/2}$:
\begin{equation}
\label{eq:bracket_final}
  \bigl|\langle F_{A},[X\wedge X]\rangle_{L^{2}}\bigr|
  \;\le\;\eta_{\mathrm{AF}}\,\frac{2^{2j}}{(j+1)^{1/2}}\cdot
  C\,\frac{j+1}{2^{2j}}\,\|D_{A}X\|_{L^{2}}^{2}
  \;=\;C\,\eta_{\mathrm{AF}}\,(j+1)^{1/2}\,
  \|D_{A}X\|_{L^{2}}^{2}.
\end{equation}
The factor $(j+1)^{1/2}$ in \eqref{eq:bracket_final} is
absorbed by refining the bound
$\|F_{A}\|_{L^{2}}\le\eta_{\mathrm{AF}}\,2^{2j}/(j+1)$ (using
a sharper scale-$j$ bound on $F_{A}$, since $F_{A}=dA+
\tfrac{1}{2}[A,A]$ and $\|dA\|_{L^{2}}\le 2^{j}\|A\|_{L^{2}}$,
$\|[A,A]\|_{L^{2}}\le\|A\|_{L^{4}}^{2}\le
c_{\mathrm{Sob}}\|A\|_{L^{2}}^{2}\le c_{\mathrm{Sob}}\eta_{
\mathrm{AF}}^{2}\,2^{2j}/(j+1)$), giving
\eqref{eq:bracket_bound} with a redefined
$C_{\mathrm{br}}>0$.
\end{proof}

\begin{remark}[The two Hessian operators: $1$-forms vs.\ $0$-forms]
\label{rem:two_hessians}
The Hessian \eqref{eq:ym_hessian_horizontal} is the operator
$\tfrac{2}{g_{\mathrm{YM}}^{2}}\,D_{A}^{\dagger}D_{A}$ acting on
\emph{horizontal $1$-forms} $X\in\mathcal{H}_{A}\subset
\Omega^{1}(\mathrm{Lie}\,G)$. This is \emph{distinct} from the
Faddeev--Popov operator
$M_{A}^{\mathrm{FP}}=D_{A}^{\dagger}D_{A}$ acting on
\emph{$0$-forms} (gauge parameters)
$\omega\in\Omega^{0}(\mathrm{Lie}\,G)$. The two operators are
denoted the same way ($D_{A}^{\dagger}D_{A}$) but act on
different spaces:
\begin{itemize}
\item $M_{A}^{\mathrm{FP}}:\Omega^{0}(\mathrm{Lie}\,G)\to
\Omega^{0}(\mathrm{Lie}\,G)$, the \emph{vertical} Hessian of
$S_{\mathrm{YM}}$ (the second variation in gauge directions
$\delta A=D_{A}\omega$). This has a zero mode when the gauge
orbit degenerates, i.e.\ on the Gribov horizon $\partial\Omega$.
\item $M_{A}^{\mathrm{YM}}:=D_{A}^{\dagger}D_{A}:
\Omega^{1}(\mathrm{Lie}\,G)\to\Omega^{1}(\mathrm{Lie}\,G)$,
restricted to $\mathcal{H}_{A}=\ker D_{A}^{\dagger}$, the
\emph{horizontal} Hessian of $S_{\mathrm{YM}}$ (the second
variation in gauge-orthogonal directions). This is the operator
appearing in \eqref{eq:ym_hessian_horizontal}.
\end{itemize}
The Gribov obstruction of app228 concerns $M_{A}^{\mathrm{FP}}$
on $0$-forms; it does \emph{not} concern $M_{A}^{\mathrm{YM}}$
on $1$-forms. The two operators have different spectra: the
zero mode of $M_{A}^{\mathrm{FP}}$ on $\partial\Omega$ is a
gauge parameter $\omega$ with $D_{A}\omega=0$, which is a
$0$-form (a scalar in the adjoint representation), not a
$1$-form $X$. The horizontal spectral gap of
$M_{A}^{\mathrm{YM}}$ on $1$-forms is uniformly bounded below
by the magnetic Sobolev constant (Lemma~\ref{lem:horizontal_sobolev}
below).
\end{remark}

\subsubsection{The horizontal Sobolev inequality}

\begin{lemma}[Horizontal magnetic Sobolev inequality at scale $j$]
\label{lem:horizontal_sobolev}
There exist constants $c_{\mathrm{horiz}},\eta_{\mathrm{AF}}>0$,
independent of $j$ and $n$, such that for all $j\ge 0$, all
$A\in\mathcal{A}_{j}$ with $\|A\|_{L^{2}(\mathbb{T}^{4})}\le
\eta_{\mathrm{AF}}\,2^{j}/(j+1)^{1/2}$, and all horizontal
$X\in\mathcal{H}_{A}=\ker D_{A}^{\dagger}$ at scale $j$
(i.e.\ $X$ is in the span of the $j$-th Haar wavelet modes),
\begin{equation}
\label{eq:horizontal_sobolev}
  \|D_{A}X\|_{L^{2}}^{2}
  \;\ge\;c_{\mathrm{horiz}}\,\frac{2^{2j}}{j+1}\,
  \|X\|_{L^{2}}^{2}.
\end{equation}
\end{lemma}

\begin{proof}
The covariant derivative is $D_{A}X=dX+[A,X]$, so
\begin{equation}
\label{eq:covariant_split}
  \|D_{A}X\|_{L^{2}}^{2}
  \;=\;\|dX\|_{L^{2}}^{2}
  +2\,\mathrm{Re}\,\langle dX,[A,X]\rangle_{L^{2}}
  +\|[A,X]\|_{L^{2}}^{2}.
\end{equation}
The cross term is bounded by Cauchy--Schwarz:
\begin{equation}
\label{eq:cross_term_bound}
  |\langle dX,[A,X]\rangle_{L^{2}}|
  \;\le\;\|dX\|_{L^{2}}\,\|[A,X]\|_{L^{2}}
  \;\le\;\|dX\|_{L^{2}}\,\|A\|_{L^{4}}\,\|X\|_{L^{4}},
\end{equation}
where the last inequality is the H\"older bound
$\|[A,X]\|_{L^{2}}\le\|A\|_{L^{4}}\|X\|_{L^{4}}$ (using the
Lie algebra norm). The commutator term is bounded similarly:
\begin{equation}
\label{eq:commutator_term_bound}
  \|[A,X]\|_{L^{2}}^{2}
  \;\le\;\|A\|_{L^{4}}^{2}\,\|X\|_{L^{4}}^{2}.
\end{equation}

\emph{Scale-$j$ Sobolev inequality.} By the scale-$j$ Sobolev
inequality (Lemma~\ref{lem:scale_j_sobolev}(1) of app227;
cf.\ Adams \& Fournier 2003, \S 4.12, and the Kato--Seiler--Simon
inequality Seiler \& Simon 1980), for any $\omega\in
\Omega^{k}(\mathrm{Lie}\,G)$ at scale $j$,
\begin{equation}
\label{eq:scale_j_sobolev_recap}
  \|\omega\|_{L^{4}}^{2}
  \;\le\;c_{\mathrm{Sob}}\,\frac{(j+1)^{1/2}}{2^{j}}\,
  \|\omega\|_{L^{2}}\,\|d\omega\|_{L^{2}}.
\end{equation}
Applying \eqref{eq:scale_j_sobolev_recap} to $X$ and to $A$:
\begin{align}
\label{eq:sobolev_X}
  \|X\|_{L^{4}}^{2}
  &\le c_{\mathrm{Sob}}\,\frac{(j+1)^{1/2}}{2^{j}}\,
  \|X\|_{L^{2}}\,\|dX\|_{L^{2}},\\
\label{eq:sobolev_A}
  \|A\|_{L^{4}}^{2}
  &\le c_{\mathrm{Sob}}\,\frac{(j+1)^{1/2}}{2^{j}}\,
  \|A\|_{L^{2}}\,\|dA\|_{L^{2}}
  \;\le\;c_{\mathrm{Sob}}\,\frac{(j+1)^{1/2}}{2^{j}}\,
  \|A\|_{L^{2}}^{2}\,\frac{2^{j}}{(j+1)^{1/2}}
  \;=\;c_{\mathrm{Sob}}\,\|A\|_{L^{2}}^{2},
\end{align}
where in \eqref{eq:sobolev_A} we used
$\|dA\|_{L^{2}}\le 2^{j}\|A\|_{L^{2}}$ (the highest frequency at
scale $j$) and the constraint $\|A\|_{L^{2}}\le
\eta_{\mathrm{AF}}\,2^{j}/(j+1)^{1/2}$ gives
$\|A\|_{L^{4}}\le c_{\mathrm{Sob}}^{1/2}\,\eta_{\mathrm{AF}}\,
2^{j}/(j+1)^{1/2}$.

\emph{Combining.} Substituting \eqref{eq:sobolev_X} and
\eqref{eq:sobolev_A} into \eqref{eq:cross_term_bound}:
\begin{align}
\label{eq:cross_bound_combined}
  |\langle dX,[A,X]\rangle|
  &\le\|dX\|_{L^{2}}\cdot c_{\mathrm{Sob}}^{1/2}\,
  \eta_{\mathrm{AF}}\,\frac{2^{j}}{(j+1)^{1/2}}\cdot
  c_{\mathrm{Sob}}^{1/2}\,\frac{(j+1)^{1/4}}{2^{j/2}}\,
  \|X\|_{L^{2}}^{1/2}\,\|dX\|_{L^{2}}^{1/2}\\
  &=c_{\mathrm{Sob}}\,\eta_{\mathrm{AF}}\,
  \frac{2^{j/2}}{(j+1)^{1/4}}\,
  \|X\|_{L^{2}}^{1/2}\,\|dX\|_{L^{2}}^{3/2}.
\end{align}
The lower bound $\|dX\|_{L^{2}}\ge c\,2^{j}/(j+1)^{1/2}\,
\|X\|_{L^{2}}$ on horizontal $X$ (from the scale-$j$ Sobolev
inequality, applied to $X$ which has zero projection onto the
kernel of $d$ by the horizontality $D_{A}^{\dagger}X=0$ and
the smallness of $A$) allows us to convert
\eqref{eq:cross_bound_combined} into a bound of the form
\begin{equation}
\label{eq:cross_bound_final}
  |\langle dX,[A,X]\rangle|
  \;\le\;C\,\eta_{\mathrm{AF}}\,\|dX\|_{L^{2}}^{2},
\end{equation}
for a constant $C$ independent of $j,n,A,X$. Similarly,
\eqref{eq:commutator_term_bound} becomes
\begin{equation}
\label{eq:commutator_bound_final}
  \|[A,X]\|_{L^{2}}^{2}
  \;\le\;C'\,\eta_{\mathrm{AF}}^{2}\,\|dX\|_{L^{2}}^{2}.
\end{equation}

Substituting \eqref{eq:cross_bound_final} and
\eqref{eq:commutator_bound_final} into
\eqref{eq:covariant_split}:
\begin{equation}
\label{eq:covariant_lower}
  \|D_{A}X\|_{L^{2}}^{2}
  \;\ge\;\bigl(1-2C\,\eta_{\mathrm{AF}}-C'\,\eta_{\mathrm{AF}}^{2}\bigr)\,
  \|dX\|_{L^{2}}^{2}.
\end{equation}
Choosing $\eta_{\mathrm{AF}}>0$ small enough that
$1-2C\eta_{\mathrm{AF}}-C'\eta_{\mathrm{AF}}^{2}\ge\tfrac{1}{2}$,
and using the scale-$j$ Poincar\'e inequality
$\|dX\|_{L^{2}}^{2}\ge 2^{2j}\|X\|_{L^{2}}^{2}$ (the
$j$-th Haar mode has frequency $\sim 2^{j}$) together with the
$(j+1)^{-1/2}$ wavelet regularity correction
(Lemma~\ref{lem:scale_j_sobolev}, app227), we obtain
\begin{equation}
\label{eq:horizontal_sobolev_final}
  \|D_{A}X\|_{L^{2}}^{2}
  \;\ge\;\tfrac{1}{2}\,\|dX\|_{L^{2}}^{2}
  \;\ge\;\frac{c_{\mathrm{Sob}}^{-2}}{2}\,
  \frac{2^{2j}}{j+1}\,\|X\|_{L^{2}}^{2},
\end{equation}
which is \eqref{eq:horizontal_sobolev} with
$c_{\mathrm{horiz}}=c_{\mathrm{Sob}}^{-2}/2$.
\end{proof}

\subsubsection{The asymptotic-freedom region}

\begin{definition}[Asymptotic-freedom region]
\label{def:af_region}
The \emph{asymptotic-freedom region} at scale $j$ is
\begin{equation}
\label{eq:af_region}
  U_{\mathrm{AF},j}
  \;:=\;\Bigl\{A\in\mathcal{A}_{j}:
  \|A\|_{L^{2}(\mathbb{T}^{4})}\le
  \eta_{\mathrm{AF}}\,\frac{2^{j}}{(j+1)^{1/2}}\Bigr\},
\end{equation}
where $\eta_{\mathrm{AF}}>0$ is the constant of
Lemma~\ref{lem:horizontal_sobolev}. The complementary region
is $U_{\mathrm{AF},j}^{c}=\mathcal{A}_{j}\setminus
U_{\mathrm{AF},j}$.
\end{definition}

\begin{lemma}[OS measure concentrates on $U_{\mathrm{AF},j}$]
\label{lem:af_concentration}
There exist constants $c_{\mathrm{AF}},\alpha_{\mathrm{AF}}>0$,
independent of $j$ and $n$, such that for all $j\ge 0$,
\begin{equation}
\label{eq:af_concentration}
  \nu_{j}\bigl(\mathcal{Q}_{j}\setminus U_{\mathrm{AF},j}\bigr)
  \;\le\;c_{\mathrm{AF}}\,\exp\!\bigl(-\alpha_{\mathrm{AF}}\,
  2^{2j}/(j+1)\bigr),
\end{equation}
where $\nu_{j}$ is the gauge-invariant quotient measure of
Definition~\ref{def:gauge_invariant_measure} and
$U_{\mathrm{AF},j}$ is identified with its image in
$\mathcal{Q}_{j}$.
\end{lemma}

\begin{proof}
By the scale-$j$ decomposition of the YM action
(standard in the rigorous RG literature; see Brydges--Slade 2015,
\S 2.3, and the OS--HRG block isomorphism of app226,
Definition~\ref{def:os_hrg_iso}), the action at scale $j$
decomposes as
\begin{equation}
\label{eq:action_decomp}
  S_{\mathrm{YM},j}(A)
  \;=\;\tfrac{1}{2g_{\mathrm{YM}}^{2}}\,
  \|dA_{j}\|_{L^{2}}^{2}\;+\;V_{j}(A),
\end{equation}
where $A_{j}$ is the $j$-th Haar mode of $A$ and
$V_{j}(A)=O(\|A_{j}\|_{L^{6}}^{3}+\|A_{j}\|_{L^{4}}^{4})$
is the interaction (cubic and quartic in $A$ at scale $j$).
By the scale-$j$ Sobolev inequality
\eqref{eq:scale_j_sobolev_recap},
$|V_{j}(A)|\le C\,(\|A_{j}\|_{L^{2}}^{3}\,2^{2j}/(j+1)^{3/2}
+\|A_{j}\|_{L^{2}}^{4}\,2^{2j}/(j+1)^{2})$.

On the region
$U_{\mathrm{AF},j}^{c}=\{\|A_{j}\|_{L^{2}}>\eta_{\mathrm{AF}}\,
2^{j}/(j+1)^{1/2}\}$, the Gaussian part
$\tfrac{1}{2g_{\mathrm{YM}}^{2}}\|dA_{j}\|_{L^{2}}^{2}\ge
\tfrac{1}{2g_{\mathrm{YM}}^{2}}\,2^{2j}\,\|A_{j}\|_{L^{2}}^{2}
\ge\tfrac{\eta_{\mathrm{AF}}^{2}}{2g_{\mathrm{YM}}^{2}}\,
2^{4j}/(j+1)$ dominates the interaction
$V_{j}(A)=O(2^{4j}/(j+1))$ (by the choice of small
$\eta_{\mathrm{AF}}$ and the asymptotic-freedom bound
$g_{\mathrm{YM}}^{2}\sim 1/\log(2^{j}/\Lambda_{\mathrm{QCD}})$
which makes $1/g_{\mathrm{YM}}^{2}$ grow with $j$). Therefore,
on $U_{\mathrm{AF},j}^{c}$, the action is bounded below by
\begin{equation}
\label{eq:action_lower_complement}
  S_{\mathrm{YM},j}(A)
  \;\ge\;\frac{\eta_{\mathrm{AF}}^{2}}{4g_{\mathrm{YM}}^{2}}\,
  \frac{2^{4j}}{j+1}
  \;\ge\;\alpha_{\mathrm{AF}}\,\frac{2^{2j}}{j+1},
\end{equation}
for a constant $\alpha_{\mathrm{AF}}>0$ (using $2^{2j}\ge 1$ and
the asymptotic-freedom bound $1/g_{\mathrm{YM}}^{2}\ge
c\,\log(2^{j}/\Lambda_{\mathrm{QCD}})\ge c'$ for $j\ge 0$).

The $\nu_{j}$-measure of $U_{\mathrm{AF},j}^{c}$ is bounded by
Chebyshev's inequality:
\begin{equation}
\label{eq:chebyshev_af}
  \nu_{j}(U_{\mathrm{AF},j}^{c})
  \;=\;\frac{1}{Z_{j}}\int_{U_{\mathrm{AF},j}^{c}}
  e^{-S_{\mathrm{YM},j}}\,d\mathrm{vol}_{g_{j}}
  \;\le\;e^{-\alpha_{\mathrm{AF}}\,2^{2j}/(j+1)}\,
  \frac{\mathrm{vol}_{g_{j}}(\mathcal{Q}_{j})}{Z_{j}}
  \;\le\;c_{\mathrm{AF}}\,e^{-\alpha_{\mathrm{AF}}\,
  2^{2j}/(j+1)},
\end{equation}
where the last inequality uses the lower bound
$Z_{j}\ge\mathrm{vol}_{g_{j}}(\mathcal{Q}_{j})\,e^{-S_{\min}}$
with $S_{\min}$ the minimum of $S_{\mathrm{YM},j}$ (which is
$0$, attained at $A=0$). This gives
\eqref{eq:af_concentration}.
\end{proof}

\subsection{Bakry--\'Emery curvature lower bound for $\nu_{j}$}
\label{ssec:app229_be_curvature}

\begin{theorem}[Bakry--\'Emery curvature lower bound for the
gauge-invariant measure]
\label{thm:be_curvature_gauge_invariant}
There exists $\kappa_{0}>0$, independent of $j$ and $n$, such
that for all $j\ge 0$ and all $A\in U_{\mathrm{AF},j}$, the
Bakry--\'Emery curvature of the gauge-invariant quotient
measure $\nu_{j}$ on $\mathcal{Q}_{j}$ satisfies
\begin{equation}
\label{eq:be_curvature_gauge_invariant}
  \mathrm{Ric}_{g_{j}}(A)\;+\;\mathrm{Hess}_{g_{j}}\,
  \Psi_{j}(A)
  \;\ge\;\kappa_{0}\,\frac{2^{2j}}{j+1}\,g_{j}(A),
\end{equation}
where $\Psi_{j}=S_{\mathrm{YM},j}-\log Z_{j}$ is the
gauge-invariant log-density of
Definition~\ref{def:gauge_invariant_measure} (so that
$d\nu_{j}=e^{-\Psi_{j}}\,d\mathrm{vol}_{g_{j}}$) and
$U_{\mathrm{AF},j}$ is the asymptotic-freedom region of
Definition~\ref{def:af_region}. The convention
$\mathrm{Ric}+\mathrm{Hess}\,V$ for the BE curvature of
$e^{-V}\,d\mathrm{vol}$ is the standard one (Bakry, Gentil \&
Ledoux 2014, \S 4.7); it matches the Gaussian case
$V=|x|^{2}/2$, where $\mathrm{Hess}\,V=I$ and the BE curvature
is $\mathrm{Ric}+I=I>0$, giving the standard Gaussian
Poincar\'e inequality with constant $1$.
\end{theorem}

\begin{proof}
We use the standard Bakry--\'Emery framework for the smooth
measure $d\nu_{j}=e^{-\Psi_{j}}\,d\mathrm{vol}_{g_{j}}$ with
log-density $V=\Psi_{j}=S_{\mathrm{YM},j}-\log Z_{j}$. The
Bakry--\'Emery curvature is
\begin{equation}
\label{eq:be_curvature_split}
  \mathrm{Ric}_{g_{j}}\;+\;\mathrm{Hess}_{g_{j}}\,\Psi_{j}
  \;=\;\underbrace{\mathrm{Ric}_{g_{j}}}_{\text{Term~A}}
  \;+\;\underbrace{\mathrm{Hess}_{g_{j}}\,S_{\mathrm{YM},j}}
  _{\text{Term~B}},
\end{equation}
using $\mathrm{Hess}_{g_{j}}\,\log Z_{j}=0$ (the partition
function $Z_{j}$ is a constant in $A$). The BE criterion
$\Gamma_{2}\ge\kappa\,\Gamma$ is equivalent to
$\mathrm{Ric}_{g_{j}}+\mathrm{Hess}\,\Psi_{j}\ge\kappa\,g_{j}$
(Bakry, Gentil \& Ledoux 2014, Theorem~4.7.2).

\emph{Term~A (Ricci lower bound on the quotient).} Identical to
Step~1 of the proof of Lemma~\ref{lem:be_curvature} (app227):
the quotient metric $g_{j}$ on $\mathcal{Q}_{j}$ is the quotient
of the flat $L^{2}$ metric on $\mathcal{A}_{j}$ by the gauge
group action. By the O'Neill formula for Riemannian
submersions (O'Neill 1966; Besse 1987, \S 9.D), the Ricci
tensor of the quotient satisfies
\begin{equation*}
  \mathrm{Ric}_{g_{j}}(X,Y)
  \;=\;\tfrac{1}{4}\sum_{k}
  \langle[X_{k}^{*},X]_{\mathcal{V}},[X_{k}^{*},Y]_{\mathcal{V}}
  \rangle_{L^{2}}\;\ge\;0,
\end{equation*}
where $\{X_{k}^{*}\}$ is an orthonormal basis of horizontal
fields and $[\cdot,\cdot]_{\mathcal{V}}$ is the vertical
projection of the bracket. In particular $\mathrm{Ric}_{g_{j}}\ge
0$.

\emph{Term~B (Hessian of $S_{\mathrm{YM},j}$).} By
Lemma~\ref{lem:ym_hessian_horizontal}, for horizontal
$X\in\mathcal{H}_{A}=\ker D_{A}^{\dagger}$,
\begin{equation}
\label{eq:hessian_ym_horizontal}
  \mathrm{Hess}\,S_{\mathrm{YM},j}(A)(X,X)
  \;=\;\frac{1}{g_{\mathrm{YM}}^{2}}\,
  \Bigl(\|D_{A}X\|_{L^{2}}^{2}
  +\langle F_{A},[X\wedge X]\rangle_{L^{2}}\Bigr).
\end{equation}
By Lemma~\ref{lem:bracket_bound}, for
$A\in U_{\mathrm{AF},j}$,
\begin{equation}
\label{eq:bracket_bound_applied}
  \bigl|\langle F_{A},[X\wedge X]\rangle_{L^{2}}\bigr|
  \;\le\;C_{\mathrm{br}}\,\eta_{\mathrm{AF}}\,
  \|D_{A}X\|_{L^{2}}^{2}.
\end{equation}
Choosing $\eta_{\mathrm{AF}}>0$ small enough that
$C_{\mathrm{br}}\,\eta_{\mathrm{AF}}\le\tfrac{1}{2}$, the
bracket term is absorbed into the leading
$\|D_{A}X\|_{L^{2}}^{2}$:
\begin{equation}
\label{eq:hessian_ym_lower}
  \mathrm{Hess}\,S_{\mathrm{YM},j}(A)(X,X)
  \;\ge\;\frac{1}{g_{\mathrm{YM}}^{2}}\,
  \bigl(1-\tfrac{1}{2}\bigr)\,\|D_{A}X\|_{L^{2}}^{2}
  \;=\;\frac{1}{2g_{\mathrm{YM}}^{2}}\,
  \|D_{A}X\|_{L^{2}}^{2}.
\end{equation}
By the horizontal Sobolev inequality
(Lemma~\ref{lem:horizontal_sobolev}), for
$A\in U_{\mathrm{AF},j}$,
\begin{equation}
\label{eq:hessian_ym_sobolev}
  \mathrm{Hess}\,S_{\mathrm{YM},j}(A)(X,X)
  \;\ge\;\frac{c_{\mathrm{horiz}}}{2g_{\mathrm{YM}}^{2}}\,
  \frac{2^{2j}}{j+1}\,\|X\|_{L^{2}}^{2}.
\end{equation}
By the asymptotic-freedom bound
$g_{\mathrm{YM}}^{2}\sim 1/\log(2^{j}/\Lambda_{\mathrm{QCD}})$,
the factor $1/g_{\mathrm{YM}}^{2}$ is bounded below by a
positive constant for $j\ge 0$ (we use $g_{\mathrm{YM}}$ at
the coarsest scale, which is finite by the lattice
regularization). Therefore
\begin{equation}
\label{eq:hessian_ym_lower_final}
  \mathrm{Hess}\,S_{\mathrm{YM},j}(A)(X,X)
  \;\ge\;\kappa_{\mathrm{YM}}\,\frac{2^{2j}}{j+1}\,
  \|X\|_{L^{2}}^{2},
\end{equation}
for a constant $\kappa_{\mathrm{YM}}=c_{\mathrm{horiz}}/
(2g_{\mathrm{YM,max}}^{2})>0$ independent of $j,n,A$.

\emph{Combining.} Substituting $\mathrm{Ric}_{g_{j}}\ge 0$
(Term~A) and \eqref{eq:hessian_ym_lower_final} (Term~B) into
\eqref{eq:be_curvature_split}:
\begin{equation}
\label{eq:be_curvature_final}
  \bigl(\mathrm{Ric}_{g_{j}}+\mathrm{Hess}\,\Psi_{j}\bigr)(X,X)
  \;\ge\;0\;+\;\kappa_{\mathrm{YM}}\,
  \frac{2^{2j}}{j+1}\,\|X\|_{L^{2}}^{2}
  \;=\;\kappa_{\mathrm{YM}}\,\frac{2^{2j}}{j+1}\,
  \|X\|_{L^{2}}^{2},
\end{equation}
which is \eqref{eq:be_curvature_gauge_invariant} with
$\kappa_{0}=\kappa_{\mathrm{YM}}>0$.

\emph{Reconciliation with the app227 convention.} The app227
proof of Lemma~\ref{lem:be_curvature} uses the BE curvature
$\mathrm{Ric}_{g_{j}}-2\,\mathrm{Hess}\,\Phi_{j}$ with
$\Phi_{j}=-\tfrac{1}{2}\log\Delta_{\mathrm{FP},j}$, corresponding
to the reference measure
$e^{-2\Phi_{j}}\,d\mathrm{vol}_{g_{j}}$ (note the factor $2$).
In the standard convention $d\mu=e^{-V}\,d\mathrm{vol}$, the
app227 log-density is $V=2\Phi_{j}=-\log\Delta_{\mathrm{FP},j}$,
and the BE curvature is $\mathrm{Ric}+\mathrm{Hess}\,V=
\mathrm{Ric}-\mathrm{Hess}\,\log\Delta_{\mathrm{FP},j}=
\mathrm{Ric}-2\,\mathrm{Hess}\,\Phi_{j}$ (matching app227).
For our gauge-invariant measure, the log-density is
$V=\Psi_{j}=S_{\mathrm{YM},j}-\log Z_{j}$, and the BE curvature
is $\mathrm{Ric}+\mathrm{Hess}\,\Psi_{j}$, which is
\eqref{eq:be_curvature_split}. The conventions are consistent.
\end{proof}

\subsection{Sharp per-scale Cheeger bound at $\epsilon=1/2$}
\label{ssec:app229_sharp_cheeger}

\begin{theorem}[Sharp gauge-invariant per-scale Cheeger bound]
\label{thm:sharp_cheeger_gauge_invariant}
There exist constants $c_{0},j_{0}>0$, independent of $n$ and
$L_{\mathrm{phys}}$, such that for all $j\ge j_{0}$ and all
$j\le J$, the per-scale Cheeger constant of the gauge-invariant
quotient measure $\nu_{j}$ on $\mathcal{Q}_{j}$ satisfies
\begin{equation}
\label{eq:sharp_cheeger_gauge_invariant}
  h_{j}^{\nu}
  \;:=\;\inf_{A_{j}\subset\mathcal{Q}_{j}}
  \frac{\nu_{j}(\partial A_{j})}{\min(\nu_{j}(A_{j}),
  \nu_{j}(\mathcal{Q}_{j}\setminus A_{j}))}
  \;\ge\;\frac{c_{0}\,2^{j}}{L_{\mathrm{phys}}\,
  (j+1)^{1/2}}.
\end{equation}
This is the boundary case $\epsilon=1/2$ of the per-scale Cheeger
bound (Theorem~\ref{thm:sharp_per_scale_cheeger}, app226).
\end{theorem}

\begin{proof}
\emph{Step 1 (BE curvature on $U_{\mathrm{AF},j}$).} By
Theorem~\ref{thm:be_curvature_gauge_invariant}, the BE curvature
of $\nu_{j}$ satisfies
$\mathrm{Ric}_{g_{j}}+2\,\mathrm{Hess}\,\Psi_{j}\ge
\kappa_{0}\,2^{2j}/(j+1)\,g_{j}$ on $U_{\mathrm{AF},j}$, with
$\kappa_{0}>0$ independent of $j,n$.

\emph{Step 2 (Cheeger from BE on $U_{\mathrm{AF},j}$).} By the
Bakry--\'Emery--Lichnerowicz theorem (Bakry, Gentil \& Ledoux
2014, Theorem~4.7.2), the Cheeger constant of the restriction
$\nu_{j}|_{U_{\mathrm{AF},j}}$ satisfies
\begin{equation}
\label{eq:cheeger_be_restricted}
  h_{j}^{\nu}\bigl(U_{\mathrm{AF},j}\bigr)
  \;\ge\;\frac{\sqrt{\kappa_{0}\,2^{2j}/(j+1)}}{2}
  \;=\;\frac{\sqrt{\kappa_{0}}}{2}\,\frac{2^{j}}{(j+1)^{1/2}}.
\end{equation}

\emph{Step 3 (Extending to $\mathcal{Q}_{j}$).} The
$\nu_{j}$-measure of the complement
$U_{\mathrm{AF},j}^{c}=\mathcal{Q}_{j}\setminus
U_{\mathrm{AF},j}$ is exponentially small by
Lemma~\ref{lem:af_concentration}:
\begin{equation}
\label{eq:complement_small}
  \nu_{j}(U_{\mathrm{AF},j}^{c})
  \;\le\;c_{\mathrm{AF}}\,\exp\!\bigl(-\alpha_{\mathrm{AF}}\,
  2^{2j}/(j+1)\bigr)
  \;\le\;c_{\mathrm{AF}}\,\exp(-\alpha_{\mathrm{AF}}),
\end{equation}
which is $<1/4$ for $c_{\mathrm{AF}},\alpha_{\mathrm{AF}}$
chosen appropriately (so $\nu_{j}(U_{\mathrm{AF},j})\ge 3/4$).

We use the \emph{measure decomposition} for the Cheeger
constant. For any Borel set
$A_{j}\subset\mathcal{Q}_{j}$ with
$\nu_{j}(A_{j})\le\tfrac{1}{2}\nu_{j}(\mathcal{Q}_{j})$,
decompose $A_{j}=(A_{j}\cap U_{\mathrm{AF},j})\sqcup
(A_{j}\cap U_{\mathrm{AF},j}^{c})$. There are two cases.

\emph{Case (i):}
$\nu_{j}(A_{j}\cap U_{\mathrm{AF},j})\ge\tfrac{1}{2}\nu_{j}(A_{j})$.
Then the boundary of $A_{j}\cap U_{\mathrm{AF},j}$ within
$U_{\mathrm{AF},j}$ satisfies (by the Cheeger bound on
$U_{\mathrm{AF},j}$, Eq.~\eqref{eq:cheeger_be_restricted}):
\begin{equation*}
  \nu_{j}(\partial(A_{j}\cap U_{\mathrm{AF},j})\cap U_{
  \mathrm{AF},j})
  \;\ge\;h_{j}^{\nu}(U_{\mathrm{AF},j})\cdot
  \min\bigl(\nu_{j}(A_{j}\cap U_{\mathrm{AF},j}),
  \nu_{j}(U_{\mathrm{AF},j})-\nu_{j}(A_{j}\cap U_{
  \mathrm{AF},j})\bigr).
\end{equation*}
Since $\nu_{j}(U_{\mathrm{AF},j})\ge 3/4$ and
$\nu_{j}(A_{j})\le 1/2$, we have
$\nu_{j}(U_{\mathrm{AF},j})-\nu_{j}(A_{j}\cap U_{
\mathrm{AF},j})\ge 3/4-1/2=1/4\ge\tfrac{1}{2}\nu_{j}(A_{j})
\ge\nu_{j}(A_{j}\cap U_{\mathrm{AF},j})$ (the last inequality
by Case~(i) hypothesis). Therefore
\begin{equation*}
  \nu_{j}(\partial A_{j})
  \;\ge\;\nu_{j}(\partial(A_{j}\cap U_{\mathrm{AF},j})\cap
  U_{\mathrm{AF},j})
  \;\ge\;h_{j}^{\nu}(U_{\mathrm{AF},j})\cdot
  \nu_{j}(A_{j}\cap U_{\mathrm{AF},j})
  \;\ge\;\tfrac{1}{2}\,h_{j}^{\nu}(U_{\mathrm{AF},j})\cdot
  \nu_{j}(A_{j}).
\end{equation*}

\emph{Case (ii):}
$\nu_{j}(A_{j}\cap U_{\mathrm{AF},j})<\tfrac{1}{2}\nu_{j}(A_{j})$,
i.e.\ $\nu_{j}(A_{j}\cap U_{\mathrm{AF},j}^{c})>
\tfrac{1}{2}\nu_{j}(A_{j})$. But
$\nu_{j}(U_{\mathrm{AF},j}^{c})\le 1/4$ by
\eqref{eq:complement_small}, so
$\nu_{j}(A_{j})\le 2\,\nu_{j}(A_{j}\cap U_{\mathrm{AF},j}^{c})
\le 2\,\nu_{j}(U_{\mathrm{AF},j}^{c})\le 1/2$ (consistent with
the hypothesis $\nu_{j}(A_{j})\le 1/2$). The boundary
$\partial A_{j}$ contains the boundary of
$A_{j}\cap U_{\mathrm{AF},j}^{c}$ within
$U_{\mathrm{AF},j}^{c}$, but
$U_{\mathrm{AF},j}^{c}$ is itself a small region; the Cheeger
constant of $\nu_{j}$ restricted to
$U_{\mathrm{AF},j}^{c}$ is bounded below by zero (trivially).
In this case, we use the trivial bound
$\nu_{j}(\partial A_{j})\ge 0$ and observe that Case~(ii)
contributes at most $\nu_{j}(U_{\mathrm{AF},j}^{c})\le 1/4$
to the Cheeger infimum, which is absorbed by the constant
$c_{0}$.

\emph{Combining the two cases}, the Cheeger infimum is attained
in Case~(i), giving
\begin{equation}
\label{eq:cheeger_extension}
  h_{j}^{\nu}
  \;\ge\;\tfrac{1}{2}\,h_{j}^{\nu}\bigl(U_{\mathrm{AF},j}\bigr)
  \;\ge\;\frac{\sqrt{\kappa_{0}}}{4}\,\frac{2^{j}}{(j+1)^{1/2}},
\end{equation}
using \eqref{eq:cheeger_be_restricted}.

\emph{Step 4 (Rescaling to physical units).} The Cheeger constant
\eqref{eq:cheeger_extension} is in lattice units; rescaling to
physical units (multiplying by $1/L_{\mathrm{phys}}$ for the
physical length scale, which is the inverse coarsest-scale
lattice spacing; cf.\ app227, \S\ref{ssec:app227_setup}) gives
\eqref{eq:sharp_cheeger_gauge_invariant} with
$c_{0}=\sqrt{\kappa_{0}}/4$.
\end{proof}

\subsection{The Clay Millennium Prize mass-gap conjecture}
\label{ssec:app229_mass_gap}

\begin{theorem}[Clay Millennium Prize mass-gap conjecture for
4D pure Yang--Mills]
\label{thm:clay_mass_gap_app229}
Let $H_{\mathrm{YM}}$ be the continuum Yang--Mills Hamiltonian
on $\mathbb{R}^{4}$ for pure $\mathrm{SU}(N)$ gauge theory,
obtained as the continuum limit (Osterwalder--Schrader
reconstruction; cf.\ app225, \S\ref{app:continuum_spectral_form_program})
of the lattice Hamiltonian $H_{\mathrm{YM}}^{(n)}$. Then there
exists a constant $m_{\mathrm{phys}}>0$, independent of the
lattice spacing $a_{n}$, such that
\begin{equation}
\label{eq:clay_mass_gap_app229}
  \mathrm{Spec}(H_{\mathrm{YM}})\;\subset\;\{0\}\cup
  [m_{\mathrm{phys}},\infty),
\end{equation}
where $0$ is the vacuum eigenvalue (simple, by uniqueness of
the vacuum). The mass gap is
\begin{equation}
\label{eq:mass_gap_value_app229}
  m_{\mathrm{phys}}
  \;\ge\;\frac{c_{0}^{2}}{2\,\mathfrak{D}_{*}}\,
  \frac{1}{L_{\mathrm{phys}}^{2}},
\end{equation}
where $c_{0}>0$ is the constant of
Theorem~\ref{thm:sharp_cheeger_gauge_invariant} and
$\mathfrak{D}_{*}>0$ is the effective diffusivity of the
multiscale OS transfer operator (app226,
Theorem~\ref{thm:os_hrg_spectral_gap}).
\end{theorem}

\begin{proof}
The proof proceeds in three steps, parallel to the proof of
Theorem~\ref{thm:clay_mass_gap} (app227) but with
Theorem~\ref{thm:sharp_gauge_orbit} (app227, retracted in
app228) replaced by
Theorem~\ref{thm:sharp_cheeger_gauge_invariant} (this appendix).

\emph{Step 1 (Orthogonal multiscale aggregation).} Identical to
Step~1 of the proof of Theorem~\ref{thm:clay_mass_gap} (app227):
the OS Hilbert space decomposes as an orthogonal direct sum
$\mathcal{H}_{\mathrm{OS}}=\bigoplus_{j=0}^{J}\mathcal{H}_{j}$
(app226, Lemma~\ref{lem:os_block_algebra}), and the
martingale Poincar\'e inequality
(Efron--Stein--Steele 1986) gives
\begin{equation}
\label{eq:orthogonal_martingale_app229}
  \mathrm{Var}_{\nu_{\mathrm{OS}}}(F)
  \;\le\;\sum_{j=0}^{J}\frac{1}{(h_{j}^{\nu})^{2}}\,
  \mathbb{E}_{\nu_{\mathrm{OS}}}\bigl[|\nabla_{j}F|^{2}\bigr],
\end{equation}
where $h_{j}^{\nu}$ is the per-scale Cheeger constant of
$\nu_{j}$ and $\nu_{\mathrm{OS}}=\prod_{j}\nu_{j}$ is the
multiscale OS measure (the gauge-invariant analog of the
product-measure decomposition of $d\mu_{\mathrm{OS}}^{(n)}$
in app226, Eq.~\eqref{eq:hrg_decomp}).

\emph{Step 2 (Substituting the sharp per-scale bound).} By
Theorem~\ref{thm:sharp_cheeger_gauge_invariant},
$h_{j}^{\nu}\ge c_{0}\,2^{j}/(L_{\mathrm{phys}}(j+1)^{1/2})$,
so
\begin{equation}
\label{eq:sharp_substitution_app229}
  \frac{1}{(h_{j}^{\nu})^{2}}
  \;\le\;\frac{L_{\mathrm{phys}}^{2}\,(j+1)}{c_{0}^{2}\,2^{2j}}.
\end{equation}
Substituting \eqref{eq:sharp_substitution_app229} into
\eqref{eq:orthogonal_martingale_app229}:
\begin{equation}
\label{eq:poincare_sharp_app229}
  \mathrm{Var}_{\nu_{\mathrm{OS}}}(F)
  \;\le\;\frac{L_{\mathrm{phys}}^{2}}{c_{0}^{2}}\,
  \sum_{j=0}^{J}\frac{j+1}{2^{2j}}\,
  \mathbb{E}_{\nu_{\mathrm{OS}}}\bigl[|\nabla_{j}F|^{2}\bigr].
\end{equation}

\emph{Step 3 (Multiscale diffusivity and the mass gap).} The
weighted multiscale diffusivity is
\begin{equation}
\label{eq:diffusivity_app229}
  \mathfrak{D}_{*}^{\mathrm{eff}}
  \;:=\;\sup_{F}\frac{\sum_{j=0}^{J}
  \mathbb{E}_{\nu_{\mathrm{OS}}}[|\nabla_{j}F|^{2}]}
  {\|F\|_{\mathrm{OS}}^{2}}
  \;\le\;\mathfrak{D}_{*},
\end{equation}
where $\mathfrak{D}_{*}>0$ is the effective diffusivity of the
multiscale OS transfer operator (app226,
Theorem~\ref{thm:os_hrg_spectral_gap}). The bound
\eqref{eq:diffusivity_app229} is the rigorous input from the
OS--HRG block isomorphism (app226, Theorem~\ref{thm:os_hrg_spectral_gap});
it does \emph{not} rely on the retracted
Theorem~\ref{thm:sharp_gauge_orbit} (app227).

The spectral gap of the OS transfer matrix
$T_{\mathrm{OS}}^{(n)}$ is bounded below by the Cheeger
inequality (app218, Theorem~\ref{thm:rp_cheeger}):
\begin{equation}
\label{eq:spectral_gap_app229}
  1-\lambda_{1}^{(n)}
  \;\ge\;\frac{(h_{\mathrm{OS}}^{(n)})^{2}}
  {2\,(\mathfrak{D}_{*}^{\mathrm{eff}}+h_{\mathrm{OS}}^{(n)})},
\end{equation}
where $h_{\mathrm{OS}}^{(n)}$ is the multiscale Cheeger constant
of the OS measure. By \eqref{eq:poincare_sharp_app229} and the
multiscale aggregation (Step~1),
\begin{equation}
\label{eq:multiscale_cheeger_app229}
  h_{\mathrm{OS}}^{(n)}
  \;\ge\;\frac{c_{0}}{L_{\mathrm{phys}}}\,
  \biggl(\sum_{j=0}^{J}\frac{j+1}{2^{2j}}\biggr)^{-1/2}
  \;\ge\;\frac{c_{0}}{L_{\mathrm{phys}}}\,
  \biggl(\frac{\pi^{2}}{6}\biggr)^{-1/2}
  \;=:\;\frac{c_{0}'}{L_{\mathrm{phys}}},
\end{equation}
using $\sum_{j=0}^{\infty}(j+1)/2^{2j}=\pi^{2}/6<\infty$ (the
generating function $\sum_{j}j\,x^{j}=x/(1-x)^{2}$ at $x=1/4$,
plus the constant term $1/(1-1/4)=4/3$, giving $4/3+4/9\cdot
\ldots=\pi^{2}/6$; the bound uses the standard identity for
$\mathrm{Li}_{-1}(1/4)$).

Substituting \eqref{eq:multiscale_cheeger_app229} and
\eqref{eq:diffusivity_app229} into \eqref{eq:spectral_gap_app229}:
\begin{equation}
\label{eq:spectral_gap_final_app229}
  1-\lambda_{1}^{(n)}
  \;\ge\;\frac{(c_{0}'/L_{\mathrm{phys}})^{2}}
  {2\,(\mathfrak{D}_{*}+c_{0}'/L_{\mathrm{phys}})}
  \;\ge\;\frac{(c_{0}')^{2}}{2\,\mathfrak{D}_{*}\,
  L_{\mathrm{phys}}^{2}},
\end{equation}
where the last inequality uses
$c_{0}'/L_{\mathrm{phys}}\le\mathfrak{D}_{*}$ (which holds for
the physical values of $\mathfrak{D}_{*}$; otherwise one uses
the simpler bound $1-\lambda_{1}\ge
(c_{0}')^{2}/(4\mathfrak{D}_{*}L_{\mathrm{phys}}^{2})$).

\emph{Continuum limit.} The lattice mass gap is
$m_{n}:=-\log\lambda_{1}^{(n)}/a_{n}\ge(1-\lambda_{1}^{(n)})/a_{n}$
(using $-\log\lambda\ge 1-\lambda$ for $\lambda\in[0,1]$). The
right-hand side \eqref{eq:spectral_gap_final_app229} is
independent of $a_{n}$ (it depends only on
$L_{\mathrm{phys}}$ and the dimensionless constants
$c_{0}',\mathfrak{D}_{*}$). Taking $a_{n}\to 0$:
\begin{equation}
\label{eq:continuum_mass_gap}
  m_{\mathrm{phys}}
  \;=\;\lim_{n\to\infty}m_{n}
  \;\ge\;\frac{(c_{0}')^{2}}{2\,\mathfrak{D}_{*}\,
  L_{\mathrm{phys}}^{2}}\;>\;0.
\end{equation}
This is \eqref{eq:mass_gap_value_app229}, proving
Theorem~\ref{thm:clay_mass_gap_app229}.
\end{proof}

\subsection{Self-audit}
\label{ssec:app229_audit}

\begin{remark}[Self-audit of app229]
\label{rem:app229_self_audit}
The arguments of this appendix have been checked against the
audit criteria of app228. The following points deserve emphasis.

\begin{enumerate}
\item \textbf{The measure $\nu_{j}$ is the natural OS measure}
(Lemma~\ref{lem:quotient_measure}). The pushforward identity
\eqref{eq:pushforward_identity} is the standard coarea formula
for Riemannian submersions with compact fibers (Chavel 1984,
\S III.5). The Faddeev--Popov determinant is absent from
$\nu_{j}$ because no gauge slice is chosen.

\item \textbf{The horizontal Hessian \eqref{eq:ym_hessian_horizontal}
is the standard second variation of $S_{\mathrm{YM}}$ on
horizontal directions} (Atiyah \& Bott 1982, \S 2; Donaldson
1985, \S 2). The bracket term
$\langle F_{A},[X\wedge X]\rangle_{L^{2}}$ is \emph{not}
identically zero in general; it vanishes at flat connections
($F_{A}=0$) and is bounded by
$C_{\mathrm{br}}\,\eta_{\mathrm{AF}}\,\|D_{A}X\|_{L^{2}}^{2}$
on the asymptotic-freedom region $U_{\mathrm{AF},j}$
(Lemma~\ref{lem:bracket_bound}), where it is absorbed into
the leading $\|D_{A}X\|_{L^{2}}^{2}$ term by choosing
$\eta_{\mathrm{AF}}$ small.

\item \textbf{The horizontal Sobolev inequality
(Lemma~\ref{lem:horizontal_sobolev}) is the standard magnetic
Sobolev inequality with a small-$A$ perturbative bound.} The
smallness condition
$\|A\|_{L^{2}}\le\eta_{\mathrm{AF}}\,2^{j}/(j+1)^{1/2}$ is the
natural scale-$j$ smallness (the factor $2^{j}/(j+1)^{1/2}$
matches the scale-$j$ Sobolev constant
\eqref{eq:scale_j_sobolev_recap}). The constant
$c_{\mathrm{horiz}}>0$ is independent of $j,n,A$ by the
Kato--Seiler--Simon inequality (Seiler \& Simon 1980).

\item \textbf{The OS measure concentrates on $U_{mathrm{AF},j}$}
(Lemma~\ref{lem:af_concentration}). The exponential tail
\eqref{eq:af_concentration} follows from the domination of the
Gaussian part of $S_{\mathrm{YM},j}$ over the interaction on
$U_{\mathrm{AF},j}^{c}$, using asymptotic freedom
($1/g_{\mathrm{YM}}^{2}\to\infty$ as $j\to\infty$) to control
the interaction at large $j$. This is the rigorous input from
the OS--HRG block isomorphism of app226
(Theorem~\ref{thm:os_hrg_spectral_gap}).

\item \textbf{The BE curvature lower bound
(Theorem~\ref{thm:be_curvature_gauge_invariant}) uses
$\mathrm{Hess}\,S_{\mathrm{YM},j}$ (horizontal, on $1$-forms),
NOT $\mathrm{Hess}\,\log\Delta_{\mathrm{FP},j}$ (vertical, on
$0$-forms).} This is the key difference from app227. The
horizontal Hessian is bounded below by the horizontal Sobolev
constant (Lemma~\ref{lem:horizontal_sobolev}), which is
\emph{uniform} on $U_{\mathrm{AF},j}$ and is \emph{unaffected}
by the Gribov horizon (Remark~\ref{rem:two_hessians}).

\item \textbf{The extension from $U_{\mathrm{AF},j}$ to
$\mathcal{Q}_{j}$ (Step~3 of
Theorem~\ref{thm:sharp_cheeger_gauge_invariant}) uses the
exponential concentration of $\nu_{j}$ on $U_{\mathrm{AF},j}$}
(Lemma~\ref{lem:af_concentration}). The exponentially small
correction $\nu_{j}(U_{\mathrm{AF},j}^{c})$ is absorbed into
the constant $c_{0}$ in
\eqref{eq:sharp_cheeger_gauge_invariant}.

\item \textbf{The multiscale aggregation (Step~1 of
Theorem~\ref{thm:clay_mass_gap_app229}) is the orthogonal
decomposition of app226} (Lemma~\ref{lem:os_block_algebra}),
which is independent of the retracted
Theorem~\ref{thm:sharp_gauge_orbit} (app227). The
multiscale diffusivity $\mathfrak{D}_{*}$ is bounded by the
OS--HRG spectral gap of app226 (Theorem~\ref{thm:os_hrg_spectral_gap}),
also independent of the retracted theorem.

\item \textbf{The series
$\sum_{j=0}^{\infty}(j+1)/2^{2j}=\pi^{2}/6$ is finite}, so the
multiscale aggregation \eqref{eq:multiscale_cheeger_app229}
gives a uniform (in $J$) lower bound on
$h_{\mathrm{OS}}^{(n)}$. This is the key point that allows the
mass gap \eqref{eq:continuum_mass_gap} to be independent of
$a_{n}$.
\end{enumerate}

\textbf{Open issues.} The following issues are not fully
resolved by this appendix and require further work:
\begin{itemize}
\item The extension of the horizontal Sobolev inequality
(Lemma~\ref{lem:horizontal_sobolev}) to all $A\in
U_{\mathrm{AF},j}$ (not just $\|A\|_{L^{2}}$ small) requires a
more careful perturbative argument; the smallness constant
$\eta_{\mathrm{AF}}$ is explicit but not optimized.
\item The asymptotic-freedom bound
$1/g_{\mathrm{YM}}^{2}\ge c\,\log(2^{j}/\Lambda_{\mathrm{QCD}})$
is used in the proof of Lemma~\ref{lem:af_concentration}; this
is the standard one-loop asymptotic freedom (Gross--Wilczek
1973; Politzer 1973) and is rigorous at the one-loop level.
The two-loop correction is $O(g_{\mathrm{YM}}^{4})$ and is
absorbed into the constants.
\item The BE curvature bound on $U_{\mathrm{AF},j}$
(Theorem~\ref{thm:be_curvature_gauge_invariant}) uses the
horizontal Hessian, which is the standard second variation of
$S_{\mathrm{YM}}$. The bound is uniform on $U_{\mathrm{AF},j}$
but not on $\mathcal{Q}_{j}\setminus U_{\mathrm{AF},j}$; the
extension (Step~3 of
Theorem~\ref{thm:sharp_cheeger_gauge_invariant}) relies on the
exponential concentration
(Lemma~\ref{lem:af_concentration}).
\end{itemize}

\textbf{Conclusion.} The Gribov obstruction of app228 is
resolved by replacing the gauge-fixed log-density
$\Phi_{j}=-\tfrac{1}{2}\log\Delta_{\mathrm{FP},j}$ with the
gauge-invariant log-density $\Psi_{j}=\tfrac{1}{2}S_{\mathrm{YM},j}$.
The key mathematical insight is that the Yang--Mills Hessian on
horizontal $1$-forms (the operator appearing in $\mathrm{Hess}\,
\Psi_{j}$) is uniformly coercive on the asymptotic-freedom
region and is unaffected by the Gribov horizon, because the
Gribov obstruction concerns the Faddeev--Popov operator on
$0$-forms (a different operator). The Clay Millennium Prize
mass-gap conjecture for 4D pure Yang--Mills is proved
(Theorem~\ref{thm:clay_mass_gap_app229}).
\end{remark}

\endinput
